%%%%%%%%%%%%%%%%%%%%%%%%%%%%%%%%%%%%%%%%%%%%%%%%%%%%%%%%%%%%
% 9.5 DIFFERENCE EQUATIONS
% The Composition of Advanced Mathematics
%%%%%%%%%%%%%%%%%%%%%%%%%%%%%%%%%%%%%%%%%%%%%%%%%%%%%%%%%%%%

\section{Difference Equations}
\label{sec:difference-equations}

\prerequisite{
Sequences, recurrence relations, algebra, functions, matrices,
eigenvalues and eigenvectors, ordinary differential equations,
dynamical systems, and basic transform methods.
}

\learningobjectives{
By the end of this section, the reader should be able to:
\begin{itemize}
    \item define a difference equation;
    \item distinguish difference equations from differential equations;
    \item determine the order of a recurrence or difference equation;
    \item distinguish linear and nonlinear difference equations;
    \item distinguish homogeneous and nonhomogeneous equations;
    \item identify autonomous and nonautonomous difference equations;
    \item solve first-order linear difference equations;
    \item solve geometric recurrences;
    \item understand equilibrium and fixed-point solutions;
    \item analyze fixed-point stability;
    \item solve higher-order constant-coefficient linear recurrences;
    \item derive and use characteristic equations;
    \item handle distinct, repeated, and complex characteristic roots;
    \item solve nonhomogeneous recurrences;
    \item use undetermined coefficients for recurrence equations;
    \item understand resonance in discrete systems;
    \item use generating functions conceptually;
    \item understand the \(z\)-transform as a solution method;
    \item formulate systems of difference equations;
    \item solve linear systems using matrix powers;
    \item analyze discrete-time stability using eigenvalues;
    \item understand Markov-type linear recurrences conceptually;
    \item model population growth and financial accumulation;
    \item analyze discrete logistic growth;
    \item distinguish continuous and discrete stability criteria;
    \item understand period-two and higher periodic orbits;
    \item recognize bifurcations in discrete maps;
    \item understand numerical discretization of differential equations;
    \item recognize finite-difference approximations;
    \item connect difference equations with dynamical systems,
          numerical analysis, probability, economics, and computation.
\end{itemize}
}

%%%%%%%%%%%%%%%%%%%%%%%%%%%%%%%%%%%%%%%%%%%%%%%%%%%%%%%%%%%%
\subsection{What Is a Difference Equation?}
\label{subsec:what-is-difference-equation}
%%%%%%%%%%%%%%%%%%%%%%%%%%%%%%%%%%%%%%%%%%%%%%%%%%%%%%%%%%%%

A difference equation relates values of a sequence at different discrete
times or indices.

A first-order example is

\[
\boxed{
x_{n+1}=2x_n.
}
\]

A second-order example is

\[
\boxed{
x_{n+2}-3x_{n+1}+2x_n=0.
}
\]

Instead of describing continuous change through derivatives, difference
equations describe discrete change through recurrence.

%%%%%%%%%%%%%%%%%%%%%%%%%%%%%%%%%%%%%%%%%%%%%%%%%%%%%%%%%%%%
\subsection{Difference Equations Versus Differential Equations}
\label{subsec:difference-versus-differential}
%%%%%%%%%%%%%%%%%%%%%%%%%%%%%%%%%%%%%%%%%%%%%%%%%%%%%%%%%%%%

A differential equation evolves continuously:

\[
\boxed{
\frac{dx}{dt}=f(x,t).
}
\]

A difference equation evolves in discrete steps:

\[
\boxed{
x_{n+1}=F(x_n,n).
}
\]

Thus,

\[
\boxed{
\text{continuous time}
\longleftrightarrow
\text{differential equation},
}
\]

while

\[
\boxed{
\text{discrete time}
\longleftrightarrow
\text{difference equation}.
}
\]

%%%%%%%%%%%%%%%%%%%%%%%%%%%%%%%%%%%%%%%%%%%%%%%%%%%%%%%%%%%%
\subsection{Forward Difference Operator}
\label{subsec:forward-difference-operator}
%%%%%%%%%%%%%%%%%%%%%%%%%%%%%%%%%%%%%%%%%%%%%%%%%%%%%%%%%%%%

\begin{definitionbox}{Forward Difference Operator}
For a sequence \(x_n\), the forward difference operator is

\[
\boxed{
\Delta x_n
=
x_{n+1}-x_n.
}
\]
\end{definitionbox}

This is the discrete analogue of a derivative.

For example, if

\[
x_n=n^2,
\]

then

\[
\Delta x_n
=
(n+1)^2-n^2
=
2n+1.
\]

%%%%%%%%%%%%%%%%%%%%%%%%%%%%%%%%%%%%%%%%%%%%%%%%%%%%%%%%%%%%
\subsection{Higher Differences}
\label{subsec:higher-differences}
%%%%%%%%%%%%%%%%%%%%%%%%%%%%%%%%%%%%%%%%%%%%%%%%%%%%%%%%%%%%

The second forward difference is

\[
\boxed{
\Delta^2x_n
=
\Delta(\Delta x_n).
}
\]

Thus,

\[
\Delta^2x_n
=
x_{n+2}
-
2x_{n+1}
+
x_n.
\]

Higher differences are defined recursively:

\[
\boxed{
\Delta^k
=
\Delta(\Delta^{k-1}).
}
\]

%%%%%%%%%%%%%%%%%%%%%%%%%%%%%%%%%%%%%%%%%%%%%%%%%%%%%%%%%%%%
\subsection{Order of a Difference Equation}
\label{subsec:difference-equation-order}
%%%%%%%%%%%%%%%%%%%%%%%%%%%%%%%%%%%%%%%%%%%%%%%%%%%%%%%%%%%%

The order is determined by the largest separation between sequence
indices needed to compute the recurrence.

For example,

\[
\boxed{
x_{n+2}
-
5x_{n+1}
+
6x_n
=
0
}
\]

is second order.

A second-order recurrence generally requires two initial values.

%%%%%%%%%%%%%%%%%%%%%%%%%%%%%%%%%%%%%%%%%%%%%%%%%%%%%%%%%%%%
\subsection{General Recurrence Form}
\label{subsec:general-difference-equation-form}
%%%%%%%%%%%%%%%%%%%%%%%%%%%%%%%%%%%%%%%%%%%%%%%%%%%%%%%%%%%%

An order-\(k\) difference equation may be written

\[
\boxed{
F
(
n,
x_n,
x_{n+1},
\ldots,
x_{n+k}
)
=
0.
}
\]

If it can be solved explicitly for the newest state,

\[
\boxed{
x_{n+k}
=
G
(
n,
x_n,\ldots,x_{n+k-1}
).
}
\]

%%%%%%%%%%%%%%%%%%%%%%%%%%%%%%%%%%%%%%%%%%%%%%%%%%%%%%%%%%%%
\subsection{Initial Conditions}
\label{subsec:difference-equation-initial-conditions}
%%%%%%%%%%%%%%%%%%%%%%%%%%%%%%%%%%%%%%%%%%%%%%%%%%%%%%%%%%%%

A first-order recurrence requires one initial value:

\[
\boxed{
x_0.
}
\]

A second-order recurrence generally requires

\[
\boxed{
x_0,
\qquad
x_1.
}
\]

An order-\(k\) recurrence typically requires \(k\) starting values.

%%%%%%%%%%%%%%%%%%%%%%%%%%%%%%%%%%%%%%%%%%%%%%%%%%%%%%%%%%%%
\subsection{Linear Difference Equations}
\label{subsec:linear-difference-equations}
%%%%%%%%%%%%%%%%%%%%%%%%%%%%%%%%%%%%%%%%%%%%%%%%%%%%%%%%%%%%

A linear order-\(k\) difference equation has the form

\[
\boxed{
a_k(n)x_{n+k}
+
a_{k-1}(n)x_{n+k-1}
+
\cdots
+
a_0(n)x_n
=
g_n.
}
\]

The sequence terms appear only to the first power and are not multiplied
together.

%%%%%%%%%%%%%%%%%%%%%%%%%%%%%%%%%%%%%%%%%%%%%%%%%%%%%%%%%%%%
\subsection{Constant-Coefficient Linear Recurrences}
\label{subsec:constant-coefficient-recurrences}
%%%%%%%%%%%%%%%%%%%%%%%%%%%%%%%%%%%%%%%%%%%%%%%%%%%%%%%%%%%%

A particularly important case is

\[
\boxed{
a_kx_{n+k}
+
a_{k-1}x_{n+k-1}
+
\cdots
+
a_0x_n
=
g_n,
}
\]

where the coefficients are constant.

These recurrences are discrete analogues of constant-coefficient linear
ODEs.

%%%%%%%%%%%%%%%%%%%%%%%%%%%%%%%%%%%%%%%%%%%%%%%%%%%%%%%%%%%%
\subsection{Homogeneous Difference Equations}
\label{subsec:homogeneous-difference-equations}
%%%%%%%%%%%%%%%%%%%%%%%%%%%%%%%%%%%%%%%%%%%%%%%%%%%%%%%%%%%%

A linear recurrence is homogeneous when

\[
\boxed{
g_n=0.
}
\]

For example,

\[
\boxed{
x_{n+2}
-
3x_{n+1}
+
2x_n
=
0.
}
\]

%%%%%%%%%%%%%%%%%%%%%%%%%%%%%%%%%%%%%%%%%%%%%%%%%%%%%%%%%%%%
\subsection{Nonhomogeneous Difference Equations}
\label{subsec:nonhomogeneous-difference-equations}
%%%%%%%%%%%%%%%%%%%%%%%%%%%%%%%%%%%%%%%%%%%%%%%%%%%%%%%%%%%%

A recurrence is nonhomogeneous when

\[
\boxed{
g_n\neq0.
}
\]

For example,

\[
\boxed{
x_{n+1}
-
2x_n
=
3.
}
\]

The sequence \(g_n\) acts as a forcing term.

%%%%%%%%%%%%%%%%%%%%%%%%%%%%%%%%%%%%%%%%%%%%%%%%%%%%%%%%%%%%
\subsection{Nonlinear Difference Equations}
\label{subsec:nonlinear-difference-equations}
%%%%%%%%%%%%%%%%%%%%%%%%%%%%%%%%%%%%%%%%%%%%%%%%%%%%%%%%%%%%

A recurrence is nonlinear if sequence values appear nonlinearly.

Examples include

\[
\boxed{
x_{n+1}=x_n^2+1,
}
\]

and

\[
\boxed{
x_{n+1}
=
rx_n(1-x_n).
}
\]

Nonlinear recurrences can exhibit periodicity, bifurcations, and chaos.

%%%%%%%%%%%%%%%%%%%%%%%%%%%%%%%%%%%%%%%%%%%%%%%%%%%%%%%%%%%%
\subsection{Autonomous Difference Equations}
\label{subsec:autonomous-difference-equations}
%%%%%%%%%%%%%%%%%%%%%%%%%%%%%%%%%%%%%%%%%%%%%%%%%%%%%%%%%%%%

A first-order autonomous difference equation has form

\[
\boxed{
x_{n+1}=f(x_n).
}
\]

The update rule does not depend explicitly on \(n\).

This is the standard form for a one-dimensional discrete dynamical
system.

%%%%%%%%%%%%%%%%%%%%%%%%%%%%%%%%%%%%%%%%%%%%%%%%%%%%%%%%%%%%
\subsection{Nonautonomous Difference Equations}
\label{subsec:nonautonomous-difference-equations}
%%%%%%%%%%%%%%%%%%%%%%%%%%%%%%%%%%%%%%%%%%%%%%%%%%%%%%%%%%%%

A nonautonomous recurrence depends explicitly on the index:

\[
\boxed{
x_{n+1}
=
f(n,x_n).
}
\]

For example,

\[
\boxed{
x_{n+1}
=
\frac{n+1}{n+2}x_n.
}
\]

%%%%%%%%%%%%%%%%%%%%%%%%%%%%%%%%%%%%%%%%%%%%%%%%%%%%%%%%%%%%
\subsection{First-Order Linear Difference Equations}
\label{subsec:first-order-linear-difference-equations}
%%%%%%%%%%%%%%%%%%%%%%%%%%%%%%%%%%%%%%%%%%%%%%%%%%%%%%%%%%%%

A first-order linear recurrence has form

\[
\boxed{
x_{n+1}
=
a_nx_n+b_n.
}
\]

If \(a_n=a\) and \(b_n=b\) are constant, then

\[
\boxed{
x_{n+1}
=
ax_n+b.
}
\]

%%%%%%%%%%%%%%%%%%%%%%%%%%%%%%%%%%%%%%%%%%%%%%%%%%%%%%%%%%%%
\subsection{Homogeneous First-Order Recurrence}
\label{subsec:first-order-homogeneous-recurrence}
%%%%%%%%%%%%%%%%%%%%%%%%%%%%%%%%%%%%%%%%%%%%%%%%%%%%%%%%%%%%

Consider

\[
\boxed{
x_{n+1}=ax_n.
}
\]

Repeated substitution gives

\[
x_1=ax_0,
\]

\[
x_2=a^2x_0,
\]

and generally

\[
\boxed{
x_n=a^nx_0.
}
\]

This is a geometric sequence.

%%%%%%%%%%%%%%%%%%%%%%%%%%%%%%%%%%%%%%%%%%%%%%%%%%%%%%%%%%%%
\subsection{Worked Example: Geometric Recurrence}
\label{subsec:worked-geometric-recurrence}
%%%%%%%%%%%%%%%%%%%%%%%%%%%%%%%%%%%%%%%%%%%%%%%%%%%%%%%%%%%%

\begin{workedexamplebox}{First-Order Homogeneous Recurrence}

Solve

\[
x_{n+1}=3x_n,
\qquad
x_0=2.
\]

The general solution is

\[
x_n=C3^n.
\]

Using

\[
x_0=2,
\]

we get

\[
C=2.
\]

Therefore,

\begin{answerbox}
\[
\boxed{
x_n=2\cdot3^n.
}
\]
\end{answerbox}

\end{workedexamplebox}

%%%%%%%%%%%%%%%%%%%%%%%%%%%%%%%%%%%%%%%%%%%%%%%%%%%%%%%%%%%%
\subsection{Constant-Forcing First-Order Recurrence}
\label{subsec:first-order-constant-forcing}
%%%%%%%%%%%%%%%%%%%%%%%%%%%%%%%%%%%%%%%%%%%%%%%%%%%%%%%%%%%%

Consider

\[
\boxed{
x_{n+1}=ax_n+b.
}
\]

If

\[
a\neq1,
\]

an equilibrium \(x^*\) satisfies

\[
x^*=ax^*+b.
\]

Thus,

\[
\boxed{
x^*
=
\frac{b}{1-a}.
}
\]

%%%%%%%%%%%%%%%%%%%%%%%%%%%%%%%%%%%%%%%%%%%%%%%%%%%%%%%%%%%%
\subsection{Shifting by the Equilibrium}
\label{subsec:equilibrium-shift-recurrence}
%%%%%%%%%%%%%%%%%%%%%%%%%%%%%%%%%%%%%%%%%%%%%%%%%%%%%%%%%%%%

Define

\[
\boxed{
y_n
=
x_n-x^*.
}
\]

Then

\[
x_{n+1}-x^*
=
a(x_n-x^*).
\]

Therefore,

\[
\boxed{
y_{n+1}=ay_n.
}
\]

Hence,

\[
\boxed{
y_n=a^ny_0.
}
\]

So,

\[
\boxed{
x_n
=
x^*
+
a^n(x_0-x^*).
}
\]

%%%%%%%%%%%%%%%%%%%%%%%%%%%%%%%%%%%%%%%%%%%%%%%%%%%%%%%%%%%%
\subsection{Closed Form for \(x_{n+1}=ax_n+b\)}
\label{subsec:first-order-affine-closed-form}
%%%%%%%%%%%%%%%%%%%%%%%%%%%%%%%%%%%%%%%%%%%%%%%%%%%%%%%%%%%%

For

\[
a\neq1,
\]

\[
\boxed{
x_n
=
a^nx_0
+
b
\frac{
1-a^n
}{
1-a
}.
}
\]

Equivalently,

\[
\boxed{
x_n
=
\frac{b}{1-a}
+
a^n
\left(
x_0-\frac{b}{1-a}
\right).
}
\]

%%%%%%%%%%%%%%%%%%%%%%%%%%%%%%%%%%%%%%%%%%%%%%%%%%%%%%%%%%%%
\subsection{Special Case \(a=1\)}
\label{subsec:first-order-a-equals-one}
%%%%%%%%%%%%%%%%%%%%%%%%%%%%%%%%%%%%%%%%%%%%%%%%%%%%%%%%%%%%

If

\[
x_{n+1}=x_n+b,
\]

then

\[
\boxed{
x_n=x_0+nb.
}
\]

Thus constant differences generate arithmetic sequences.

%%%%%%%%%%%%%%%%%%%%%%%%%%%%%%%%%%%%%%%%%%%%%%%%%%%%%%%%%%%%
\subsection{Worked Example: Affine Recurrence}
\label{subsec:worked-affine-recurrence}
%%%%%%%%%%%%%%%%%%%%%%%%%%%%%%%%%%%%%%%%%%%%%%%%%%%%%%%%%%%%

\begin{workedexamplebox}{Convergence to an Equilibrium}

Solve

\[
x_{n+1}
=
\frac12x_n+3,
\qquad
x_0=10.
\]

The equilibrium satisfies

\[
x^*
=
\frac12x^*+3.
\]

Thus

\[
x^*=6.
\]

Therefore,

\[
x_n
=
6
+
\left(
\frac12
\right)^n
(10-6).
\]

Hence,

\begin{answerbox}
\[
\boxed{
x_n
=
6
+
4
\left(
\frac12
\right)^n.
}
\]

Since

\[
\left(
\frac12
\right)^n\to0,
\]

\[
\boxed{
x_n\to6.
}
\]
\end{answerbox}

\end{workedexamplebox}

%%%%%%%%%%%%%%%%%%%%%%%%%%%%%%%%%%%%%%%%%%%%%%%%%%%%%%%%%%%%
\subsection{Equilibrium Solutions}
\label{subsec:difference-equation-equilibria}
%%%%%%%%%%%%%%%%%%%%%%%%%%%%%%%%%%%%%%%%%%%%%%%%%%%%%%%%%%%%

For

\[
x_{n+1}=f(x_n),
\]

an equilibrium or fixed point satisfies

\[
\boxed{
x^*=f(x^*).
}
\]

If

\[
x_0=x^*,
\]

then

\[
x_n=x^*
\]

for every \(n\).

%%%%%%%%%%%%%%%%%%%%%%%%%%%%%%%%%%%%%%%%%%%%%%%%%%%%%%%%%%%%
\subsection{Local Stability of Fixed Points}
\label{subsec:difference-fixed-point-stability}
%%%%%%%%%%%%%%%%%%%%%%%%%%%%%%%%%%%%%%%%%%%%%%%%%%%%%%%%%%%%

Suppose

\[
x^*=f(x^*).
\]

If

\[
\boxed{
|f'(x^*)|<1,
}
\]

then \(x^*\) is locally asymptotically stable.

If

\[
\boxed{
|f'(x^*)|>1,
}
\]

then \(x^*\) is unstable.

If

\[
\boxed{
|f'(x^*)|=1,
}
\]

the derivative test is inconclusive.

%%%%%%%%%%%%%%%%%%%%%%%%%%%%%%%%%%%%%%%%%%%%%%%%%%%%%%%%%%%%
\subsection{Linear First-Order Stability}
\label{subsec:linear-first-order-difference-stability}
%%%%%%%%%%%%%%%%%%%%%%%%%%%%%%%%%%%%%%%%%%%%%%%%%%%%%%%%%%%%

For

\[
x_{n+1}=ax_n+b,
\]

the equilibrium is stable precisely when

\[
\boxed{
|a|<1.
}
\]

If

\[
0<a<1,
\]

convergence is monotone on each side.

If

\[
-1<a<0,
\]

the sequence alternates around the equilibrium while converging.

%%%%%%%%%%%%%%%%%%%%%%%%%%%%%%%%%%%%%%%%%%%%%%%%%%%%%%%%%%%%
\subsection{Higher-Order Linear Homogeneous Recurrences}
\label{subsec:higher-order-linear-homogeneous-recurrences}
%%%%%%%%%%%%%%%%%%%%%%%%%%%%%%%%%%%%%%%%%%%%%%%%%%%%%%%%%%%%

Consider

\[
\boxed{
a_kx_{n+k}
+
a_{k-1}x_{n+k-1}
+
\cdots
+
a_0x_n
=
0.
}
\]

Try a solution of the form

\[
\boxed{
x_n=r^n.
}
\]

Substitution produces the characteristic polynomial.

%%%%%%%%%%%%%%%%%%%%%%%%%%%%%%%%%%%%%%%%%%%%%%%%%%%%%%%%%%%%
\subsection{Characteristic Equation}
\label{subsec:difference-characteristic-equation}
%%%%%%%%%%%%%%%%%%%%%%%%%%%%%%%%%%%%%%%%%%%%%%%%%%%%%%%%%%%%

For

\[
a_kx_{n+k}
+
\cdots
+
a_0x_n
=
0,
\]

substituting \(x_n=r^n\) gives

\[
\boxed{
a_kr^k
+
a_{k-1}r^{k-1}
+
\cdots
+
a_0
=
0.
}
\]

This is the characteristic equation.

%%%%%%%%%%%%%%%%%%%%%%%%%%%%%%%%%%%%%%%%%%%%%%%%%%%%%%%%%%%%
\subsection{Second-Order Characteristic Equation}
\label{subsec:second-order-difference-characteristic}
%%%%%%%%%%%%%%%%%%%%%%%%%%%%%%%%%%%%%%%%%%%%%%%%%%%%%%%%%%%%

For

\[
\boxed{
x_{n+2}
+
ax_{n+1}
+
bx_n
=
0,
}
\]

the characteristic equation is

\[
\boxed{
r^2+ar+b=0.
}
\]

The form of the general solution depends on the roots.

%%%%%%%%%%%%%%%%%%%%%%%%%%%%%%%%%%%%%%%%%%%%%%%%%%%%%%%%%%%%
\subsection{Distinct Real Characteristic Roots}
\label{subsec:difference-distinct-real-roots}
%%%%%%%%%%%%%%%%%%%%%%%%%%%%%%%%%%%%%%%%%%%%%%%%%%%%%%%%%%%%

If the roots are distinct,

\[
r_1\neq r_2,
\]

then

\[
\boxed{
x_n
=
C_1r_1^n
+
C_2r_2^n.
}
\]

%%%%%%%%%%%%%%%%%%%%%%%%%%%%%%%%%%%%%%%%%%%%%%%%%%%%%%%%%%%%
\subsection{Worked Example: Distinct Roots}
\label{subsec:worked-difference-distinct-roots}
%%%%%%%%%%%%%%%%%%%%%%%%%%%%%%%%%%%%%%%%%%%%%%%%%%%%%%%%%%%%

\begin{workedexamplebox}{Second-Order Recurrence}

Solve

\[
x_{n+2}
-
5x_{n+1}
+
6x_n
=
0.
\]

The characteristic equation is

\[
r^2-5r+6=0.
\]

Factor:

\[
(r-2)(r-3)=0.
\]

Thus,

\[
r_1=2,
\qquad
r_2=3.
\]

Therefore,

\begin{answerbox}
\[
\boxed{
x_n
=
C_1 2^n
+
C_2 3^n.
}
\]
\end{answerbox}

\end{workedexamplebox}

%%%%%%%%%%%%%%%%%%%%%%%%%%%%%%%%%%%%%%%%%%%%%%%%%%%%%%%%%%%%
\subsection{Using Initial Values}
\label{subsec:difference-initial-values}
%%%%%%%%%%%%%%%%%%%%%%%%%%%%%%%%%%%%%%%%%%%%%%%%%%%%%%%%%%%%

Suppose additionally

\[
x_0=1,
\qquad
x_1=4.
\]

Then

\[
C_1+C_2=1,
\]

and

\[
2C_1+3C_2=4.
\]

Solving gives

\[
C_1=-1,
\qquad
C_2=2.
\]

Hence,

\[
\boxed{
x_n
=
-2^n
+
2\cdot3^n.
}
\]

%%%%%%%%%%%%%%%%%%%%%%%%%%%%%%%%%%%%%%%%%%%%%%%%%%%%%%%%%%%%
\subsection{Repeated Characteristic Roots}
\label{subsec:difference-repeated-roots}
%%%%%%%%%%%%%%%%%%%%%%%%%%%%%%%%%%%%%%%%%%%%%%%%%%%%%%%%%%%%

If the characteristic equation has repeated root \(r\) of multiplicity
\(m\), then the corresponding solutions are

\[
\boxed{
r^n,
\quad
nr^n,
\quad
n^2r^n,
\ldots,
n^{m-1}r^n.
}
\]

For a double root,

\[
\boxed{
x_n
=
(C_1+C_2n)r^n.
}
\]

%%%%%%%%%%%%%%%%%%%%%%%%%%%%%%%%%%%%%%%%%%%%%%%%%%%%%%%%%%%%
\subsection{Worked Example: Repeated Root}
\label{subsec:worked-difference-repeated-root}
%%%%%%%%%%%%%%%%%%%%%%%%%%%%%%%%%%%%%%%%%%%%%%%%%%%%%%%%%%%%

\begin{workedexamplebox}{Repeated Characteristic Root}

Solve

\[
x_{n+2}
-
4x_{n+1}
+
4x_n
=
0.
\]

The characteristic equation is

\[
r^2-4r+4=0,
\]

so

\[
(r-2)^2=0.
\]

Therefore,

\begin{answerbox}
\[
\boxed{
x_n
=
(C_1+C_2n)2^n.
}
\]
\end{answerbox}

\end{workedexamplebox}

%%%%%%%%%%%%%%%%%%%%%%%%%%%%%%%%%%%%%%%%%%%%%%%%%%%%%%%%%%%%
\subsection{Complex Characteristic Roots}
\label{subsec:difference-complex-roots}
%%%%%%%%%%%%%%%%%%%%%%%%%%%%%%%%%%%%%%%%%%%%%%%%%%%%%%%%%%%%

Suppose the roots are

\[
\boxed{
r
=
\rho e^{\pm i\theta}.
}
\]

Then a real-valued solution can be written

\[
\boxed{
x_n
=
\rho^n
\left[
C_1\cos(n\theta)
+
C_2\sin(n\theta)
\right].
}
\]

The magnitude \(\rho\) determines growth or decay.

%%%%%%%%%%%%%%%%%%%%%%%%%%%%%%%%%%%%%%%%%%%%%%%%%%%%%%%%%%%%
\subsection{Discrete Oscillation}
\label{subsec:discrete-oscillation}
%%%%%%%%%%%%%%%%%%%%%%%%%%%%%%%%%%%%%%%%%%%%%%%%%%%%%%%%%%%%

If

\[
\boxed{
\rho<1,
}
\]

oscillations decay.

If

\[
\boxed{
\rho=1,
}
\]

oscillation magnitude remains constant in the ideal linear recurrence.

If

\[
\boxed{
\rho>1,
}
\]

oscillation magnitude grows.

%%%%%%%%%%%%%%%%%%%%%%%%%%%%%%%%%%%%%%%%%%%%%%%%%%%%%%%%%%%%
\subsection{Worked Example: Complex Roots}
\label{subsec:worked-difference-complex-roots}
%%%%%%%%%%%%%%%%%%%%%%%%%%%%%%%%%%%%%%%%%%%%%%%%%%%%%%%%%%%%

\begin{workedexamplebox}{Oscillatory Recurrence}

Consider

\[
x_{n+2}
-
x_{n+1}
+
x_n
=
0.
\]

The characteristic equation is

\[
r^2-r+1=0.
\]

Thus,

\[
r
=
\frac12
\pm
i\frac{\sqrt3}{2}
=
e^{\pm i\pi/3}.
\]

Therefore,

\begin{answerbox}
\[
\boxed{
x_n
=
C_1\cos
\left(
\frac{n\pi}{3}
\right)
+
C_2\sin
\left(
\frac{n\pi}{3}
\right).
}
\]
\end{answerbox}

\end{workedexamplebox}

%%%%%%%%%%%%%%%%%%%%%%%%%%%%%%%%%%%%%%%%%%%%%%%%%%%%%%%%%%%%
\subsection{Superposition Principle}
\label{subsec:difference-superposition}
%%%%%%%%%%%%%%%%%%%%%%%%%%%%%%%%%%%%%%%%%%%%%%%%%%%%%%%%%%%%

For a homogeneous linear recurrence

\[
L[x_n]=0,
\]

if \(x_n^{(1)}\) and \(x_n^{(2)}\) are solutions, then

\[
\boxed{
C_1x_n^{(1)}
+
C_2x_n^{(2)}
}
\]

is also a solution.

This is the discrete analogue of superposition for linear ODEs.

%%%%%%%%%%%%%%%%%%%%%%%%%%%%%%%%%%%%%%%%%%%%%%%%%%%%%%%%%%%%
\subsection{Nonhomogeneous Linear Recurrences}
\label{subsec:nonhomogeneous-linear-recurrences}
%%%%%%%%%%%%%%%%%%%%%%%%%%%%%%%%%%%%%%%%%%%%%%%%%%%%%%%%%%%%

For

\[
\boxed{
L[x_n]=g_n,
}
\]

the general solution has form

\[
\boxed{
x_n
=
x_n^{(h)}
+
x_n^{(p)},
}
\]

where

\[
x_n^{(h)}
\]

solves the homogeneous recurrence and

\[
x_n^{(p)}
\]

is one particular solution.

%%%%%%%%%%%%%%%%%%%%%%%%%%%%%%%%%%%%%%%%%%%%%%%%%%%%%%%%%%%%
\subsection{Undetermined Coefficients for Recurrences}
\label{subsec:undetermined-coefficients-recurrences}
%%%%%%%%%%%%%%%%%%%%%%%%%%%%%%%%%%%%%%%%%%%%%%%%%%%%%%%%%%%%

For constant-coefficient recurrences, particular solutions may often be
guessed from the forcing term.

If

\[
g_n=C,
\]

try a constant.

If

\[
g_n=an+b,
\]

try a polynomial of the same degree.

If

\[
g_n=c\lambda^n,
\]

try

\[
\boxed{
x_n^{(p)}=A\lambda^n.
}
\]

%%%%%%%%%%%%%%%%%%%%%%%%%%%%%%%%%%%%%%%%%%%%%%%%%%%%%%%%%%%%
\subsection{Discrete Resonance}
\label{subsec:discrete-resonance}
%%%%%%%%%%%%%%%%%%%%%%%%%%%%%%%%%%%%%%%%%%%%%%%%%%%%%%%%%%%%

If the proposed forcing trial is already a homogeneous solution, multiply
by a sufficient power of \(n\).

If \(\lambda\) is a characteristic root of multiplicity \(m\), then for
forcing proportional to

\[
\lambda^n,
\]

try

\[
\boxed{
x_n^{(p)}
=
An^m\lambda^n.
}
\]

%%%%%%%%%%%%%%%%%%%%%%%%%%%%%%%%%%%%%%%%%%%%%%%%%%%%%%%%%%%%
\subsection{Worked Example: Nonhomogeneous Recurrence}
\label{subsec:worked-nonhomogeneous-recurrence}
%%%%%%%%%%%%%%%%%%%%%%%%%%%%%%%%%%%%%%%%%%%%%%%%%%%%%%%%%%%%

\begin{workedexamplebox}{Constant Forcing}

Solve

\[
x_{n+1}
-
2x_n
=
3.
\]

The homogeneous equation is

\[
x_{n+1}-2x_n=0,
\]

with solution

\[
x_n^{(h)}
=
C2^n.
\]

Try a constant particular solution

\[
x_n^{(p)}=A.
\]

Then

\[
A-2A=3,
\]

so

\[
A=-3.
\]

Therefore,

\begin{answerbox}
\[
\boxed{
x_n
=
C2^n-3.
}
\]
\end{answerbox}

\end{workedexamplebox}

%%%%%%%%%%%%%%%%%%%%%%%%%%%%%%%%%%%%%%%%%%%%%%%%%%%%%%%%%%%%
\subsection{Worked Example: Resonant Forcing}
\label{subsec:worked-discrete-resonance}
%%%%%%%%%%%%%%%%%%%%%%%%%%%%%%%%%%%%%%%%%%%%%%%%%%%%%%%%%%%%

\begin{workedexamplebox}{Forcing Matches a Characteristic Root}

Solve

\[
x_{n+1}
-
2x_n
=
2^n.
\]

The homogeneous solution is

\[
x_n^{(h)}
=
C2^n.
\]

A trial \(A2^n\) duplicates the homogeneous solution.

Therefore try

\[
x_n^{(p)}
=
An2^n.
\]

Then

\[
x_{n+1}^{(p)}
=
A(n+1)2^{n+1}.
\]

Substitution gives

\[
A(n+1)2^{n+1}
-
2An2^n
=
2^n.
\]

Simplify:

\[
2A2^n
=
2^n.
\]

Thus

\[
A=\frac12.
\]

Hence,

\begin{answerbox}
\[
\boxed{
x_n
=
C2^n
+
\frac12n2^n.
}
\]
\end{answerbox}

\end{workedexamplebox}

%%%%%%%%%%%%%%%%%%%%%%%%%%%%%%%%%%%%%%%%%%%%%%%%%%%%%%%%%%%%
\subsection{Variable-Coefficient First-Order Recurrences}
\label{subsec:variable-coefficient-first-order-recurrence}
%%%%%%%%%%%%%%%%%%%%%%%%%%%%%%%%%%%%%%%%%%%%%%%%%%%%%%%%%%%%

Consider

\[
\boxed{
x_{n+1}
=
a_nx_n+b_n.
}
\]

Repeated substitution gives

\[
x_n
=
\left(
\prod_{j=0}^{n-1}a_j
\right)x_0
+
\sum_{k=0}^{n-1}
\left[
b_k
\prod_{j=k+1}^{n-1}
a_j
\right],
\]

with empty products interpreted as \(1\).

This is the discrete analogue of an integrating-factor formula.

%%%%%%%%%%%%%%%%%%%%%%%%%%%%%%%%%%%%%%%%%%%%%%%%%%%%%%%%%%%%
\subsection{Discrete Integrating Factor Concept}
\label{subsec:discrete-integrating-factor}
%%%%%%%%%%%%%%%%%%%%%%%%%%%%%%%%%%%%%%%%%%%%%%%%%%%%%%%%%%%%

A recurrence such as

\[
x_{n+1}-a_nx_n=b_n
\]

can sometimes be multiplied by a sequence \(\mu_n\) so the left side
becomes a telescoping difference.

The goal is analogous to producing

\[
(\mu y)'
\]

in a first-order linear ODE.

%%%%%%%%%%%%%%%%%%%%%%%%%%%%%%%%%%%%%%%%%%%%%%%%%%%%%%%%%%%%
\subsection{Telescoping Sums}
\label{subsec:difference-telescoping}
%%%%%%%%%%%%%%%%%%%%%%%%%%%%%%%%%%%%%%%%%%%%%%%%%%%%%%%%%%%%

If a recurrence can be written

\[
\boxed{
y_{n+1}-y_n=g_n,
}
\]

then summing from \(n=0\) to \(N-1\) gives

\[
y_N-y_0
=
\sum_{n=0}^{N-1}g_n.
\]

Hence,

\[
\boxed{
y_N
=
y_0
+
\sum_{n=0}^{N-1}g_n.
}
\]

Telescoping plays the role of integration in many discrete problems.

%%%%%%%%%%%%%%%%%%%%%%%%%%%%%%%%%%%%%%%%%%%%%%%%%%%%%%%%%%%%
\subsection{Fibonacci Recurrence}
\label{subsec:fibonacci-difference-equation}
%%%%%%%%%%%%%%%%%%%%%%%%%%%%%%%%%%%%%%%%%%%%%%%%%%%%%%%%%%%%

The Fibonacci recurrence is

\[
\boxed{
F_{n+2}
=
F_{n+1}
+
F_n,
}
\]

with

\[
F_0=0,
\qquad
F_1=1.
\]

The characteristic equation is

\[
\boxed{
r^2-r-1=0.
}
\]

Its roots are

\[
\boxed{
\phi
=
\frac{
1+\sqrt5
}{
2
},
}
\]

and

\[
\boxed{
\psi
=
\frac{
1-\sqrt5
}{
2
}.
}
\]

%%%%%%%%%%%%%%%%%%%%%%%%%%%%%%%%%%%%%%%%%%%%%%%%%%%%%%%%%%%%
\subsection{Binet Formula}
\label{subsec:binet-formula}
%%%%%%%%%%%%%%%%%%%%%%%%%%%%%%%%%%%%%%%%%%%%%%%%%%%%%%%%%%%%

Using the initial conditions,

\[
\boxed{
F_n
=
\frac{
\phi^n-\psi^n
}{
\sqrt5
}.
}
\]

Since

\[
|\psi|<1,
\]

the large-\(n\) behavior is approximately

\[
\boxed{
F_n
\sim
\frac{
\phi^n
}{
\sqrt5
}.
}
\]

%%%%%%%%%%%%%%%%%%%%%%%%%%%%%%%%%%%%%%%%%%%%%%%%%%%%%%%%%%%%
\subsection{Generating Functions}
\label{subsec:generating-functions-difference-equations}
%%%%%%%%%%%%%%%%%%%%%%%%%%%%%%%%%%%%%%%%%%%%%%%%%%%%%%%%%%%%

A generating function encodes a sequence as a power series.

For

\[
x_0,x_1,x_2,\ldots,
\]

define

\[
\boxed{
G(z)
=
\sum_{n=0}^{\infty}
x_nz^n.
}
\]

Recurrences can often be transformed into algebraic equations for
\(G(z)\).

%%%%%%%%%%%%%%%%%%%%%%%%%%%%%%%%%%%%%%%%%%%%%%%%%%%%%%%%%%%%
\subsection{Generating Function Strategy}
\label{subsec:generating-function-strategy}
%%%%%%%%%%%%%%%%%%%%%%%%%%%%%%%%%%%%%%%%%%%%%%%%%%%%%%%%%%%%

A typical method is:

\begin{enumerate}
    \item multiply the recurrence by \(z^n\);
    \item sum over the allowed indices;
    \item rewrite shifted sums using \(G(z)\);
    \item solve algebraically for \(G(z)\);
    \item expand or decompose \(G(z)\);
    \item read off the sequence coefficients.
\end{enumerate}

%%%%%%%%%%%%%%%%%%%%%%%%%%%%%%%%%%%%%%%%%%%%%%%%%%%%%%%%%%%%
\subsection{Worked Example: Geometric Generating Function}
\label{subsec:worked-geometric-generating-function}
%%%%%%%%%%%%%%%%%%%%%%%%%%%%%%%%%%%%%%%%%%%%%%%%%%%%%%%%%%%%

\begin{workedexamplebox}{Generating Function for \(a^n\)}

Suppose

\[
x_n=a^n.
\]

Then

\[
G(z)
=
\sum_{n=0}^{\infty}
a^nz^n.
\]

This is a geometric series:

\[
G(z)
=
\sum_{n=0}^{\infty}
(az)^n.
\]

Therefore,

\begin{answerbox}
\[
\boxed{
G(z)
=
\frac1{
1-az
},
\qquad
|az|<1.
}
\]
\end{answerbox}

\end{workedexamplebox}

%%%%%%%%%%%%%%%%%%%%%%%%%%%%%%%%%%%%%%%%%%%%%%%%%%%%%%%%%%%%
\subsection{Generating Function for Fibonacci Numbers}
\label{subsec:fibonacci-generating-function}
%%%%%%%%%%%%%%%%%%%%%%%%%%%%%%%%%%%%%%%%%%%%%%%%%%%%%%%%%%%%

For

\[
F_{n+2}=F_{n+1}+F_n,
\]

with

\[
F_0=0,
\qquad
F_1=1,
\]

the ordinary generating function is

\[
\boxed{
G(z)
=
\frac{
z
}{
1-z-z^2
}.
}
\]

Factoring the denominator and using partial fractions leads to Binet's
formula.

%%%%%%%%%%%%%%%%%%%%%%%%%%%%%%%%%%%%%%%%%%%%%%%%%%%%%%%%%%%%
\subsection{\(z\)-Transform}
\label{subsec:z-transform-difference-equations}
%%%%%%%%%%%%%%%%%%%%%%%%%%%%%%%%%%%%%%%%%%%%%%%%%%%%%%%%%%%%

The one-sided \(z\)-transform of a sequence \(x_n\) is

\[
\boxed{
X(z)
=
\sum_{n=0}^{\infty}
x_nz^{-n}.
}
\]

It is closely related to generating functions under the substitution

\[
z\mapsto z^{-1}.
\]

%%%%%%%%%%%%%%%%%%%%%%%%%%%%%%%%%%%%%%%%%%%%%%%%%%%%%%%%%%%%
\subsection{Shift Property of the One-Sided \(z\)-Transform}
\label{subsec:z-transform-shift-property}
%%%%%%%%%%%%%%%%%%%%%%%%%%%%%%%%%%%%%%%%%%%%%%%%%%%%%%%%%%%%

If

\[
X(z)
=
\sum_{n=0}^{\infty}
x_nz^{-n},
\]

then

\[
\boxed{
\mathcal{Z}\{x_{n+1}\}
=
z
[
X(z)-x_0
].
}
\]

Similarly,

\[
\boxed{
\mathcal{Z}\{x_{n+2}\}
=
z^2
\left[
X(z)-x_0-x_1z^{-1}
\right].
}
\]

These formulas incorporate initial conditions directly.

%%%%%%%%%%%%%%%%%%%%%%%%%%%%%%%%%%%%%%%%%%%%%%%%%%%%%%%%%%%%
\subsection{Solving Recurrences with the \(z\)-Transform}
\label{subsec:solving-recurrences-z-transform}
%%%%%%%%%%%%%%%%%%%%%%%%%%%%%%%%%%%%%%%%%%%%%%%%%%%%%%%%%%%%

The \(z\)-transform method parallels the Laplace transform:

\[
\boxed{
\text{difference equation}
\longrightarrow
\text{algebraic equation in }z
\longrightarrow
\text{inverse transform}.
}
\]

It is especially useful for:

\[
\boxed{
\text{linear discrete-time systems},
}
\]

\[
\boxed{
\text{digital signal processing},
}
\]

and

\[
\boxed{
\text{control systems}.
}
\]

%%%%%%%%%%%%%%%%%%%%%%%%%%%%%%%%%%%%%%%%%%%%%%%%%%%%%%%%%%%%
\subsection{Discrete-Time Transfer Function}
\label{subsec:discrete-time-transfer-function}
%%%%%%%%%%%%%%%%%%%%%%%%%%%%%%%%%%%%%%%%%%%%%%%%%%%%%%%%%%%%

For a linear discrete-time input-output system with zero initial
conditions,

\[
\boxed{
H(z)
=
\frac{
Y(z)
}{
U(z)
}
}
\]

is the transfer function.

Its poles and zeros determine important features of system behavior.

%%%%%%%%%%%%%%%%%%%%%%%%%%%%%%%%%%%%%%%%%%%%%%%%%%%%%%%%%%%%
\subsection{Discrete-Time Pole Stability}
\label{subsec:discrete-pole-stability}
%%%%%%%%%%%%%%%%%%%%%%%%%%%%%%%%%%%%%%%%%%%%%%%%%%%%%%%%%%%%

For a finite-dimensional discrete-time linear system, asymptotic
stability requires all relevant poles to lie strictly inside the unit
circle:

\[
\boxed{
|z|<1.
}
\]

This contrasts with continuous-time Laplace stability, where poles must
lie in the left half-plane.

%%%%%%%%%%%%%%%%%%%%%%%%%%%%%%%%%%%%%%%%%%%%%%%%%%%%%%%%%%%%
\subsection{Systems of Difference Equations}
\label{subsec:systems-of-difference-equations}
%%%%%%%%%%%%%%%%%%%%%%%%%%%%%%%%%%%%%%%%%%%%%%%%%%%%%%%%%%%%

A discrete system may be written

\[
\boxed{
\mathbf{x}_{n+1}
=
\mathbf{F}(\mathbf{x}_n).
}
\]

A linear system has form

\[
\boxed{
\mathbf{x}_{n+1}
=
A\mathbf{x}_n.
}
\]

%%%%%%%%%%%%%%%%%%%%%%%%%%%%%%%%%%%%%%%%%%%%%%%%%%%%%%%%%%%%
\subsection{Solution of a Linear Discrete System}
\label{subsec:linear-discrete-system-solution}
%%%%%%%%%%%%%%%%%%%%%%%%%%%%%%%%%%%%%%%%%%%%%%%%%%%%%%%%%%%%

Starting from

\[
\mathbf{x}_0,
\]

we have

\[
\mathbf{x}_1=A\mathbf{x}_0,
\]

\[
\mathbf{x}_2=A^2\mathbf{x}_0,
\]

and generally

\[
\boxed{
\mathbf{x}_n
=
A^n\mathbf{x}_0.
}
\]

Thus matrix powers play the role that \(e^{At}\) plays for continuous
linear systems.

%%%%%%%%%%%%%%%%%%%%%%%%%%%%%%%%%%%%%%%%%%%%%%%%%%%%%%%%%%%%
\subsection{Eigenvalue Solution of a Discrete Linear System}
\label{subsec:eigenvalue-discrete-system}
%%%%%%%%%%%%%%%%%%%%%%%%%%%%%%%%%%%%%%%%%%%%%%%%%%%%%%%%%%%%

Suppose \(A\) has eigenvector \(\mathbf{v}\) with eigenvalue \(\lambda\):

\[
A\mathbf{v}
=
\lambda\mathbf{v}.
\]

Then

\[
A^n\mathbf{v}
=
\lambda^n\mathbf{v}.
\]

Therefore,

\[
\boxed{
\mathbf{x}_n
=
\lambda^n\mathbf{v}
}
\]

is a mode of the system.

%%%%%%%%%%%%%%%%%%%%%%%%%%%%%%%%%%%%%%%%%%%%%%%%%%%%%%%%%%%%
\subsection{Discrete Linear Stability}
\label{subsec:discrete-linear-stability}
%%%%%%%%%%%%%%%%%%%%%%%%%%%%%%%%%%%%%%%%%%%%%%%%%%%%%%%%%%%%

For

\[
\mathbf{x}_{n+1}=A\mathbf{x}_n,
\]

the origin is asymptotically stable if every eigenvalue satisfies

\[
\boxed{
|\lambda_j|<1.
}
\]

If any eigenvalue satisfies

\[
\boxed{
|\lambda_j|>1,
}
\]

the origin is unstable.

Eigenvalues on the unit circle require more careful analysis.

%%%%%%%%%%%%%%%%%%%%%%%%%%%%%%%%%%%%%%%%%%%%%%%%%%%%%%%%%%%%
\subsection{Spectral Radius}
\label{subsec:discrete-spectral-radius}
%%%%%%%%%%%%%%%%%%%%%%%%%%%%%%%%%%%%%%%%%%%%%%%%%%%%%%%%%%%%

The spectral radius of \(A\) is

\[
\boxed{
\rho(A)
=
\max_j
|\lambda_j|.
}
\]

For a finite-dimensional linear system,

\[
\boxed{
\rho(A)<1
}
\]

implies

\[
A^n\to0.
\]

Therefore,

\[
\boxed{
\mathbf{x}_n\to0.
}
\]

%%%%%%%%%%%%%%%%%%%%%%%%%%%%%%%%%%%%%%%%%%%%%%%%%%%%%%%%%%%%
\subsection{Worked Example: Matrix Difference Equation}
\label{subsec:worked-matrix-difference-equation}
%%%%%%%%%%%%%%%%%%%%%%%%%%%%%%%%%%%%%%%%%%%%%%%%%%%%%%%%%%%%

\begin{workedexamplebox}{Diagonal Discrete System}

Consider

\[
\mathbf{x}_{n+1}
=
\begin{pmatrix}
1/2&0\\
0&-1/3
\end{pmatrix}
\mathbf{x}_n.
\]

Then

\[
A^n
=
\begin{pmatrix}
(1/2)^n&0\\
0&(-1/3)^n
\end{pmatrix}.
\]

Hence,

\[
\begin{answerbox}
\[
\boxed{
\mathbf{x}_n
=
\begin{pmatrix}
(1/2)^n&0\\
0&(-1/3)^n
\end{pmatrix}
\mathbf{x}_0.
}
\]

Since both eigenvalues have magnitude less than \(1\),

\[
\boxed{
\mathbf{x}_n\to\mathbf{0}.
}
\]
\end{answerbox}

\end{workedexamplebox}

%%%%%%%%%%%%%%%%%%%%%%%%%%%%%%%%%%%%%%%%%%%%%%%%%%%%%%%%%%%%
\subsection{Nonhomogeneous Linear Systems}
\label{subsec:nonhomogeneous-discrete-linear-systems}
%%%%%%%%%%%%%%%%%%%%%%%%%%%%%%%%%%%%%%%%%%%%%%%%%%%%%%%%%%%%

Consider

\[
\boxed{
\mathbf{x}_{n+1}
=
A\mathbf{x}_n
+
\mathbf{b}_n.
}
\]

Repeated substitution gives

\[
\boxed{
\mathbf{x}_n
=
A^n\mathbf{x}_0
+
\sum_{k=0}^{n-1}
A^{n-1-k}
\mathbf{b}_k.
}
\]

This is the discrete variation-of-constants formula.

%%%%%%%%%%%%%%%%%%%%%%%%%%%%%%%%%%%%%%%%%%%%%%%%%%%%%%%%%%%%
\subsection{Constant Input}
\label{subsec:discrete-system-constant-input}
%%%%%%%%%%%%%%%%%%%%%%%%%%%%%%%%%%%%%%%%%%%%%%%%%%%%%%%%%%%%

If

\[
\mathbf{b}_n=\mathbf{b},
\]

then

\[
\boxed{
\mathbf{x}_n
=
A^n\mathbf{x}_0
+
\sum_{k=0}^{n-1}
A^k\mathbf{b},
}
\]

after reversing the summation index.

If \(I-A\) is invertible,

\[
\boxed{
\sum_{k=0}^{n-1}
A^k
=
(I-A^n)(I-A)^{-1}.
}
\]

%%%%%%%%%%%%%%%%%%%%%%%%%%%%%%%%%%%%%%%%%%%%%%%%%%%%%%%%%%%%
\subsection{Discrete Equilibrium of a Linear System}
\label{subsec:discrete-linear-system-equilibrium}
%%%%%%%%%%%%%%%%%%%%%%%%%%%%%%%%%%%%%%%%%%%%%%%%%%%%%%%%%%%%

For

\[
\mathbf{x}_{n+1}
=
A\mathbf{x}_n+\mathbf{b},
\]

an equilibrium satisfies

\[
\mathbf{x}^*
=
A\mathbf{x}^*
+
\mathbf{b}.
\]

Therefore,

\[
\boxed{
(I-A)\mathbf{x}^*
=
\mathbf{b}.
}
\]

If \(I-A\) is invertible,

\[
\boxed{
\mathbf{x}^*
=
(I-A)^{-1}\mathbf{b}.
}
\]

%%%%%%%%%%%%%%%%%%%%%%%%%%%%%%%%%%%%%%%%%%%%%%%%%%%%%%%%%%%%
\subsection{Population Growth Model}
\label{subsec:discrete-population-growth}
%%%%%%%%%%%%%%%%%%%%%%%%%%%%%%%%%%%%%%%%%%%%%%%%%%%%%%%%%%%%

Suppose a population grows by a fixed percentage per time step.

Then

\[
\boxed{
P_{n+1}
=
(1+r)P_n.
}
\]

Therefore,

\[
\boxed{
P_n
=
P_0(1+r)^n.
}
\]

This is the discrete analogue of exponential growth.

%%%%%%%%%%%%%%%%%%%%%%%%%%%%%%%%%%%%%%%%%%%%%%%%%%%%%%%%%%%%
\subsection{Discrete Versus Continuous Growth}
\label{subsec:discrete-versus-continuous-growth}
%%%%%%%%%%%%%%%%%%%%%%%%%%%%%%%%%%%%%%%%%%%%%%%%%%%%%%%%%%%%

Continuous growth:

\[
\boxed{
P(t)
=
P_0e^{kt}.
}
\]

Discrete growth:

\[
\boxed{
P_n
=
P_0(1+r)^n.
}
\]

The two are related by

\[
\boxed{
1+r=e^k
}
\]

when the time step is one unit and the models are matched exactly at
integer times.

%%%%%%%%%%%%%%%%%%%%%%%%%%%%%%%%%%%%%%%%%%%%%%%%%%%%%%%%%%%%
\subsection{Compound Interest}
\label{subsec:difference-compound-interest}
%%%%%%%%%%%%%%%%%%%%%%%%%%%%%%%%%%%%%%%%%%%%%%%%%%%%%%%%%%%%

If an account earns rate \(r\) per period,

\[
\boxed{
A_{n+1}
=
(1+r)A_n.
}
\]

Therefore,

\[
\boxed{
A_n
=
A_0(1+r)^n.
}
\]

If a constant deposit \(d\) is added at the end of each period,

\[
\boxed{
A_{n+1}
=
(1+r)A_n+d.
}
\]

%%%%%%%%%%%%%%%%%%%%%%%%%%%%%%%%%%%%%%%%%%%%%%%%%%%%%%%%%%%%
\subsection{Savings Recurrence}
\label{subsec:savings-recurrence}
%%%%%%%%%%%%%%%%%%%%%%%%%%%%%%%%%%%%%%%%%%%%%%%%%%%%%%%%%%%%

For

\[
A_{n+1}
=
qA_n+d,
\]

where

\[
q=1+r,
\]

the solution is

\[
\boxed{
A_n
=
q^nA_0
+
d
\frac{
q^n-1
}{
q-1
},
}
\]

provided

\[
q\neq1.
\]

This is the standard future-value structure for regular deposits.

%%%%%%%%%%%%%%%%%%%%%%%%%%%%%%%%%%%%%%%%%%%%%%%%%%%%%%%%%%%%
\subsection{Loan Balance Recurrence}
\label{subsec:loan-balance-recurrence}
%%%%%%%%%%%%%%%%%%%%%%%%%%%%%%%%%%%%%%%%%%%%%%%%%%%%%%%%%%%%

If interest is added and a payment \(p\) is then made,

\[
\boxed{
B_{n+1}
=
(1+r)B_n-p.
}
\]

This is a first-order nonhomogeneous linear difference equation.

Its equilibrium is

\[
\boxed{
B^*
=
\frac{p}{r}
}
\]

when \(r>0\), although the practical goal is usually repayment rather
than convergence to this equilibrium.

%%%%%%%%%%%%%%%%%%%%%%%%%%%%%%%%%%%%%%%%%%%%%%%%%%%%%%%%%%%%
\subsection{Discrete Logistic Model}
\label{subsec:discrete-logistic-model}
%%%%%%%%%%%%%%%%%%%%%%%%%%%%%%%%%%%%%%%%%%%%%%%%%%%%%%%%%%%%

A discrete analogue of logistic growth may be written

\[
\boxed{
P_{n+1}
=
P_n
+
rP_n
\left(
1-\frac{P_n}{K}
\right).
}
\]

This differs from the classical logistic map

\[
x_{n+1}
=
Rx_n(1-x_n),
\]

although they are related through scaling and parameter choices.

%%%%%%%%%%%%%%%%%%%%%%%%%%%%%%%%%%%%%%%%%%%%%%%%%%%%%%%%%%%%
\subsection{Equilibria of the Discrete Logistic Growth Model}
\label{subsec:discrete-logistic-growth-equilibria}
%%%%%%%%%%%%%%%%%%%%%%%%%%%%%%%%%%%%%%%%%%%%%%%%%%%%%%%%%%%%

For

\[
P_{n+1}
=
P_n
+
rP_n
\left(
1-\frac{P_n}{K}
\right),
\]

an equilibrium satisfies

\[
rP
\left(
1-\frac{P}{K}
\right)=0.
\]

Thus,

\[
\boxed{
P^*=0,
}
\]

and

\[
\boxed{
P^*=K.
}
\]

%%%%%%%%%%%%%%%%%%%%%%%%%%%%%%%%%%%%%%%%%%%%%%%%%%%%%%%%%%%%
\subsection{Classical Logistic Map}
\label{subsec:difference-logistic-map}
%%%%%%%%%%%%%%%%%%%%%%%%%%%%%%%%%%%%%%%%%%%%%%%%%%%%%%%%%%%%

The classical logistic map is

\[
\boxed{
x_{n+1}
=
Rx_n(1-x_n).
}
\]

It is a nonlinear first-order difference equation.

Its fixed points are

\[
\boxed{
x^*=0
}
\]

and

\[
\boxed{
x^*
=
1-\frac1R
}
\]

when \(R\neq0\).

%%%%%%%%%%%%%%%%%%%%%%%%%%%%%%%%%%%%%%%%%%%%%%%%%%%%%%%%%%%%
\subsection{Logistic Map Stability}
\label{subsec:difference-logistic-map-stability}
%%%%%%%%%%%%%%%%%%%%%%%%%%%%%%%%%%%%%%%%%%%%%%%%%%%%%%%%%%%%

For

\[
f(x)=Rx(1-x),
\]

\[
f'(x)=R(1-2x).
\]

At \(x=0\),

\[
\boxed{
f'(0)=R.
}
\]

At the nonzero fixed point,

\[
\boxed{
f'(x^*)=2-R.
}
\]

Therefore the nonzero fixed point is stable for

\[
\boxed{
1<R<3.
}
\]

%%%%%%%%%%%%%%%%%%%%%%%%%%%%%%%%%%%%%%%%%%%%%%%%%%%%%%%%%%%%
\subsection{Period Doubling in Difference Equations}
\label{subsec:difference-period-doubling}
%%%%%%%%%%%%%%%%%%%%%%%%%%%%%%%%%%%%%%%%%%%%%%%%%%%%%%%%%%%%

As a parameter varies, a stable fixed point can lose stability and a
stable period-two orbit can appear.

This is a period-doubling bifurcation.

Repeated period doubling can produce

\[
\boxed{
2,\ 4,\ 8,\ 16,\ldots
}
\]

cycles before chaotic behavior emerges.

%%%%%%%%%%%%%%%%%%%%%%%%%%%%%%%%%%%%%%%%%%%%%%%%%%%%%%%%%%%%
\subsection{Periodic Points}
\label{subsec:difference-periodic-points}
%%%%%%%%%%%%%%%%%%%%%%%%%%%%%%%%%%%%%%%%%%%%%%%%%%%%%%%%%%%%

A point \(x\) has period \(k\) if

\[
\boxed{
f^k(x)=x
}
\]

and \(k\) is the smallest positive integer with that property.

Fixed points are period-one points.

A period-two point satisfies

\[
\boxed{
f^2(x)=x
}
\]

but

\[
\boxed{
f(x)\neq x.
}
\]

%%%%%%%%%%%%%%%%%%%%%%%%%%%%%%%%%%%%%%%%%%%%%%%%%%%%%%%%%%%%
\subsection{Stability of a Periodic Orbit}
\label{subsec:difference-periodic-orbit-stability}
%%%%%%%%%%%%%%%%%%%%%%%%%%%%%%%%%%%%%%%%%%%%%%%%%%%%%%%%%%%%

For a period-\(k\) orbit

\[
x_1,\ldots,x_k,
\]

its multiplier is

\[
\boxed{
\mu
=
\prod_{j=1}^{k}
f'(x_j).
}
\]

The orbit is locally attracting if

\[
\boxed{
|\mu|<1.
}
\]

It is repelling if

\[
\boxed{
|\mu|>1.
}
\]

%%%%%%%%%%%%%%%%%%%%%%%%%%%%%%%%%%%%%%%%%%%%%%%%%%%%%%%%%%%%
\subsection{Cobweb Diagrams}
\label{subsec:difference-cobweb-diagrams}
%%%%%%%%%%%%%%%%%%%%%%%%%%%%%%%%%%%%%%%%%%%%%%%%%%%%%%%%%%%%

For

\[
x_{n+1}=f(x_n),
\]

a cobweb diagram compares

\[
\boxed{
y=f(x)
}
\]

with

\[
\boxed{
y=x.
}
\]

Starting from \(x_0\), alternate vertically to \(y=f(x)\) and horizontally
to \(y=x\).

This visualizes iteration and fixed-point stability.

%%%%%%%%%%%%%%%%%%%%%%%%%%%%%%%%%%%%%%%%%%%%%%%%%%%%%%%%%%%%
\subsection{Monotone and Oscillatory Convergence}
\label{subsec:discrete-monotone-oscillatory-convergence}
%%%%%%%%%%%%%%%%%%%%%%%%%%%%%%%%%%%%%%%%%%%%%%%%%%%%%%%%%%%%

Near an attracting fixed point,

\[
e_{n+1}
\approx
f'(x^*)e_n.
\]

If

\[
0<f'(x^*)<1,
\]

the errors tend to retain their sign, producing monotone-type convergence.

If

\[
-1<f'(x^*)<0,
\]

the errors change sign, producing alternating convergence.

%%%%%%%%%%%%%%%%%%%%%%%%%%%%%%%%%%%%%%%%%%%%%%%%%%%%%%%%%%%%
\subsection{Discrete Bifurcations}
\label{subsec:discrete-bifurcations}
%%%%%%%%%%%%%%%%%%%%%%%%%%%%%%%%%%%%%%%%%%%%%%%%%%%%%%%%%%%%

As parameters vary, discrete maps can undergo:

\[
\boxed{
\text{saddle-node bifurcations},
}
\]

\[
\boxed{
\text{transcritical bifurcations},
}
\]

\[
\boxed{
\text{pitchfork bifurcations},
}
\]

and

\[
\boxed{
\text{period-doubling bifurcations}.
}
\]

The stability boundary for a scalar fixed point occurs when

\[
\boxed{
|f'(x^*)|=1.
}
\]

%%%%%%%%%%%%%%%%%%%%%%%%%%%%%%%%%%%%%%%%%%%%%%%%%%%%%%%%%%%%
\subsection{Difference Equations from Numerical ODE Methods}
\label{subsec:difference-equations-from-odes}
%%%%%%%%%%%%%%%%%%%%%%%%%%%%%%%%%%%%%%%%%%%%%%%%%%%%%%%%%%%%

Numerical methods convert differential equations into difference
equations.

For

\[
y'=f(t,y),
\]

Euler's method gives

\[
\boxed{
y_{n+1}
=
y_n
+
h
f(t_n,y_n).
}
\]

Thus numerical integration replaces continuous dynamics by a discrete
recurrence.

%%%%%%%%%%%%%%%%%%%%%%%%%%%%%%%%%%%%%%%%%%%%%%%%%%%%%%%%%%%%
\subsection{Euler Discretization of Exponential Decay}
\label{subsec:euler-discrete-decay}
%%%%%%%%%%%%%%%%%%%%%%%%%%%%%%%%%%%%%%%%%%%%%%%%%%%%%%%%%%%%

Consider

\[
y'=-ky,
\qquad
k>0.
\]

Euler's method gives

\[
y_{n+1}
=
y_n
-
hky_n.
\]

Therefore,

\[
\boxed{
y_{n+1}
=
(1-hk)y_n.
}
\]

Hence,

\[
\boxed{
y_n
=
(1-hk)^ny_0.
}
\]

%%%%%%%%%%%%%%%%%%%%%%%%%%%%%%%%%%%%%%%%%%%%%%%%%%%%%%%%%%%%
\subsection{Numerical Stability of Euler Decay}
\label{subsec:euler-decay-stability}
%%%%%%%%%%%%%%%%%%%%%%%%%%%%%%%%%%%%%%%%%%%%%%%%%%%%%%%%%%%%

The exact differential equation is asymptotically stable.

The Euler recurrence is stable only if

\[
\boxed{
|1-hk|<1.
}
\]

Thus,

\[
-1<1-hk<1.
\]

For \(h>0\) and \(k>0\),

\[
\boxed{
0<hk<2.
}
\]

Therefore the numerical step size can create artificial instability.

%%%%%%%%%%%%%%%%%%%%%%%%%%%%%%%%%%%%%%%%%%%%%%%%%%%%%%%%%%%%
\subsection{Finite Differences}
\label{subsec:finite-differences-difference-equations}
%%%%%%%%%%%%%%%%%%%%%%%%%%%%%%%%%%%%%%%%%%%%%%%%%%%%%%%%%%%%

Derivatives can be approximated by differences.

The forward difference is

\[
\boxed{
f'(x)
\approx
\frac{
f(x+h)-f(x)
}{
h
}.
}
\]

The centered difference is

\[
\boxed{
f'(x)
\approx
\frac{
f(x+h)-f(x-h)
}{
2h
}.
}
\]

A second derivative can be approximated by

\[
\boxed{
f''(x)
\approx
\frac{
f(x+h)-2f(x)+f(x-h)
}{
h^2
}.
}
\]

%%%%%%%%%%%%%%%%%%%%%%%%%%%%%%%%%%%%%%%%%%%%%%%%%%%%%%%%%%%%
\subsection{Discrete Boundary-Value Problem}
\label{subsec:discrete-boundary-value-problem}
%%%%%%%%%%%%%%%%%%%%%%%%%%%%%%%%%%%%%%%%%%%%%%%%%%%%%%%%%%%%

Suppose

\[
y''(x)=f(x).
\]

Using a grid

\[
x_i=ih,
\]

the second derivative gives

\[
\boxed{
\frac{
y_{i+1}-2y_i+y_{i-1}
}{
h^2
}
=
f_i.
}
\]

Thus a differential boundary-value problem becomes a system of linear
difference equations.

%%%%%%%%%%%%%%%%%%%%%%%%%%%%%%%%%%%%%%%%%%%%%%%%%%%%%%%%%%%%
\subsection{Markov-Type Recurrences}
\label{subsec:markov-type-recurrences}
%%%%%%%%%%%%%%%%%%%%%%%%%%%%%%%%%%%%%%%%%%%%%%%%%%%%%%%%%%%%

Probability distributions can evolve by matrix recurrence.

If

\[
\mathbf{p}_n
\]

is a probability vector and \(P\) is a transition matrix, one common
convention is

\[
\boxed{
\mathbf{p}_{n+1}
=
P^T\mathbf{p}_n.
}
\]

Alternatively, row-vector conventions use

\[
\boxed{
\mathbf{p}_{n+1}
=
\mathbf{p}_nP.
}
\]

The convention must be stated consistently.

%%%%%%%%%%%%%%%%%%%%%%%%%%%%%%%%%%%%%%%%%%%%%%%%%%%%%%%%%%%%
\subsection{Stationary Distribution}
\label{subsec:stationary-distribution-difference}
%%%%%%%%%%%%%%%%%%%%%%%%%%%%%%%%%%%%%%%%%%%%%%%%%%%%%%%%%%%%

A stationary distribution satisfies

\[
\boxed{
\mathbf{p}^*
=
P^T\mathbf{p}^*
}
\]

under the column-vector convention.

Thus \(\mathbf{p}^*\) is an eigenvector corresponding to eigenvalue \(1\),
subject to

\[
\boxed{
p_i^*\geq0,
}
\]

and

\[
\boxed{
\sum_i p_i^*=1.
}
\]

%%%%%%%%%%%%%%%%%%%%%%%%%%%%%%%%%%%%%%%%%%%%%%%%%%%%%%%%%%%%
\subsection{Leslie Population Matrix Preview}
\label{subsec:leslie-matrix-preview}
%%%%%%%%%%%%%%%%%%%%%%%%%%%%%%%%%%%%%%%%%%%%%%%%%%%%%%%%%%%%

Age-structured populations can be modeled by

\[
\boxed{
\mathbf{n}_{k+1}
=
L\mathbf{n}_k,
}
\]

where \(L\) is a Leslie matrix.

The dominant eigenvalue determines the long-run population growth factor
under suitable assumptions.

Its eigenvector determines the stable age distribution.

%%%%%%%%%%%%%%%%%%%%%%%%%%%%%%%%%%%%%%%%%%%%%%%%%%%%%%%%%%%%
\subsection{Long-Term Behavior of Linear Recurrences}
\label{subsec:long-term-linear-recurrence}
%%%%%%%%%%%%%%%%%%%%%%%%%%%%%%%%%%%%%%%%%%%%%%%%%%%%%%%%%%%%

For a solution built from characteristic roots,

\[
x_n
=
\sum_j
P_j(n)r_j^n,
\]

the largest root magnitude generally dominates long-term behavior.

If a unique dominant root satisfies

\[
|r_1|>|r_j|
\]

for all \(j\neq1\), then asymptotically

\[
\boxed{
x_n
\sim
P_1(n)r_1^n,
}
\]

provided its coefficient is nonzero.

%%%%%%%%%%%%%%%%%%%%%%%%%%%%%%%%%%%%%%%%%%%%%%%%%%%%%%%%%%%%
\subsection{Discrete Green's Function Preview}
\label{subsec:discrete-greens-function-preview}
%%%%%%%%%%%%%%%%%%%%%%%%%%%%%%%%%%%%%%%%%%%%%%%%%%%%%%%%%%%%

Linear nonhomogeneous recurrences can be represented using a discrete
impulse response.

For a linear time-invariant recurrence, the response to a general input
can be expressed as a discrete convolution.

This parallels Green's functions for differential equations.

%%%%%%%%%%%%%%%%%%%%%%%%%%%%%%%%%%%%%%%%%%%%%%%%%%%%%%%%%%%%
\subsection{Discrete Convolution}
\label{subsec:discrete-convolution}
%%%%%%%%%%%%%%%%%%%%%%%%%%%%%%%%%%%%%%%%%%%%%%%%%%%%%%%%%%%%

For sequences \(f_n\) and \(g_n\), a one-sided discrete convolution is

\[
\boxed{
(f*g)_n
=
\sum_{k=0}^{n}
f_k g_{n-k}.
}
\]

Under suitable \(z\)-transform conventions,

\[
\boxed{
\mathcal{Z}\{f*g\}
=
F(z)G(z).
}
\]

%%%%%%%%%%%%%%%%%%%%%%%%%%%%%%%%%%%%%%%%%%%%%%%%%%%%%%%%%%%%
\subsection{Difference Equation Workflow}
\label{subsec:difference-equation-workflow}
%%%%%%%%%%%%%%%%%%%%%%%%%%%%%%%%%%%%%%%%%%%%%%%%%%%%%%%%%%%%

A practical workflow is:

\begin{enumerate}
    \item determine the order of the recurrence;
    \item identify the required initial values;
    \item classify the equation as linear or nonlinear;
    \item determine whether coefficients are constant;
    \item solve the homogeneous equation;
    \item construct the characteristic polynomial when appropriate;
    \item determine a particular solution for nonhomogeneous forcing;
    \item apply initial conditions;
    \item examine equilibria and fixed points;
    \item analyze stability;
    \item inspect long-term behavior;
    \item use generating functions or the \(z\)-transform when useful;
    \item use numerical iteration for nonlinear recurrences.
\end{enumerate}

%%%%%%%%%%%%%%%%%%%%%%%%%%%%%%%%%%%%%%%%%%%%%%%%%%%%%%%%%%%%
\subsection{Choosing a Solution Method}
\label{subsec:choosing-difference-solution-method}
%%%%%%%%%%%%%%%%%%%%%%%%%%%%%%%%%%%%%%%%%%%%%%%%%%%%%%%%%%%%

For

\[
x_{n+1}=ax_n+b,
\]

use:

\[
\boxed{
\text{direct iteration or equilibrium shifting}.
}
\]

For constant-coefficient homogeneous higher-order recurrences, use:

\[
\boxed{
\text{characteristic roots}.
}
\]

For simple nonhomogeneous forcing, use:

\[
\boxed{
\text{undetermined coefficients}.
}
\]

For combinatorial recurrences, consider:

\[
\boxed{
\text{generating functions}.
}
\]

For discrete linear systems, consider:

\[
\boxed{
\text{matrix powers or the }z\text{-transform}.
}
\]

%%%%%%%%%%%%%%%%%%%%%%%%%%%%%%%%%%%%%%%%%%%%%%%%%%%%%%%%%%%%
\subsection{Common Errors}
\label{subsec:difference-equation-common-errors}
%%%%%%%%%%%%%%%%%%%%%%%%%%%%%%%%%%%%%%%%%%%%%%%%%%%%%%%%%%%%

\subsubsection{Using Continuous-Time Stability Rules}

For

\[
x_{n+1}=\lambda x_n,
\]

stability requires

\[
\boxed{
|\lambda|<1.
}
\]

It is not enough for

\[
\lambda<0.
\]

%%%%%%%%%%%%%%%%%%%%%%%%%%%%%%%%%%%%%%%%%%%%%%%%%%%%%%%%%%%%
\subsubsection{Forgetting Initial Conditions}

An order-\(k\) recurrence generally requires \(k\) initial values.

Without enough initial data, the constants in the general solution cannot
be determined.

%%%%%%%%%%%%%%%%%%%%%%%%%%%%%%%%%%%%%%%%%%%%%%%%%%%%%%%%%%%%
\subsubsection{Using the Wrong Characteristic Trial}

For constant-coefficient recurrences, the natural trial is

\[
\boxed{
x_n=r^n.
}
\]

This is the discrete counterpart of \(e^{rt}\) for ODEs.

%%%%%%%%%%%%%%%%%%%%%%%%%%%%%%%%%%%%%%%%%%%%%%%%%%%%%%%%%%%%
\subsubsection{Forgetting Multiplicity}

A repeated characteristic root requires factors of \(n\):

\[
\boxed{
r^n,
\ nr^n,
\ n^2r^n,
\ldots
}
\]

according to multiplicity.

%%%%%%%%%%%%%%%%%%%%%%%%%%%%%%%%%%%%%%%%%%%%%%%%%%%%%%%%%%%%
\subsubsection{Ignoring Resonance}

If a particular-solution trial duplicates a homogeneous solution, multiply
by a sufficient power of \(n\).

%%%%%%%%%%%%%%%%%%%%%%%%%%%%%%%%%%%%%%%%%%%%%%%%%%%%%%%%%%%%
\subsubsection{Confusing Fixed Points with Zeros of \(f\)}

For

\[
x_{n+1}=f(x_n),
\]

fixed points satisfy

\[
\boxed{
f(x^*)=x^*,
}
\]

not generally

\[
f(x^*)=0.
\]

%%%%%%%%%%%%%%%%%%%%%%%%%%%%%%%%%%%%%%%%%%%%%%%%%%%%%%%%%%%%
\subsubsection{Ignoring Negative Stable Multipliers}

If

\[
-1<f'(x^*)<0,
\]

the fixed point can still be stable.

Convergence typically oscillates across the equilibrium.

%%%%%%%%%%%%%%%%%%%%%%%%%%%%%%%%%%%%%%%%%%%%%%%%%%%%%%%%%%%%
\subsubsection{Assuming Boundedness Implies Convergence}

A bounded sequence may oscillate periodically or chaotically.

Boundedness alone does not imply convergence to a fixed point.

%%%%%%%%%%%%%%%%%%%%%%%%%%%%%%%%%%%%%%%%%%%%%%%%%%%%%%%%%%%%
\subsubsection{Confusing Periodic Orbits with Fixed Points}

A period-two orbit contains two distinct points.

Neither point is individually fixed under one application of \(f\).

%%%%%%%%%%%%%%%%%%%%%%%%%%%%%%%%%%%%%%%%%%%%%%%%%%%%%%%%%%%%
\subsubsection{Ignoring the Unit Circle in Discrete Linear Systems}

For matrix recurrence

\[
\mathbf{x}_{n+1}=A\mathbf{x}_n,
\]

eigenvalues are compared with the unit circle, not the imaginary axis.

%%%%%%%%%%%%%%%%%%%%%%%%%%%%%%%%%%%%%%%%%%%%%%%%%%%%%%%%%%%%
\subsubsection{Confusing Generating Functions and \(z\)-Transforms}

They are closely related but use different power conventions:

\[
G(z)
=
\sum x_nz^n,
\]

while a common one-sided \(z\)-transform is

\[
X(z)
=
\sum x_nz^{-n}.
\]

%%%%%%%%%%%%%%%%%%%%%%%%%%%%%%%%%%%%%%%%%%%%%%%%%%%%%%%%%%%%
\subsubsection{Ignoring Numerical Instability in Discretization}

A stable differential equation can generate an unstable numerical
difference equation if the step size is inappropriate.

%%%%%%%%%%%%%%%%%%%%%%%%%%%%%%%%%%%%%%%%%%%%%%%%%%%%%%%%%%%%
\subsubsection{Using a Markov Matrix Convention Inconsistently}

Row-vector and column-vector conventions transpose the transition matrix.

Choose one convention and keep it throughout the calculation.

%%%%%%%%%%%%%%%%%%%%%%%%%%%%%%%%%%%%%%%%%%%%%%%%%%%%%%%%%%%%
\subsection{Applications}
\label{subsec:difference-equation-applications}
%%%%%%%%%%%%%%%%%%%%%%%%%%%%%%%%%%%%%%%%%%%%%%%%%%%%%%%%%%%%

\begin{applicationbox}

\textbf{Population Biology.}
Year-to-year population changes, age-structured populations, and
generation-based reproduction naturally use difference equations.

\medskip

\textbf{Finance.}
Compound interest, savings plans, loan balances, and periodic payments
are modeled by first-order linear recurrences.

\medskip

\textbf{Economics.}
Discrete-time growth, cobweb models, inventory adjustment, and business
cycles use recurrence equations.

\medskip

\textbf{Probability.}
Markov chains evolve probability distributions through matrix difference
equations.

\medskip

\textbf{Computer Science.}
Algorithmic complexity and recursive procedures generate recurrence
relations.

\medskip

\textbf{Engineering.}
Digital control systems and sampled-data systems are modeled in discrete
time.

\medskip

\textbf{Signal Processing.}
Digital filters satisfy linear difference equations and are analyzed
using the \(z\)-transform.

\medskip

\textbf{Numerical Analysis.}
Finite-difference and time-stepping methods convert differential equations
into discrete recurrence systems.

\end{applicationbox}

%%%%%%%%%%%%%%%%%%%%%%%%%%%%%%%%%%%%%%%%%%%%%%%%%%%%%%%%%%%%
\subsection{Exercises}
\label{subsec:difference-equation-exercises}
%%%%%%%%%%%%%%%%%%%%%%%%%%%%%%%%%%%%%%%%%%%%%%%%%%%%%%%%%%%%

\begin{exercise}[1]
Define a difference equation.
\end{exercise}

\begin{exercise}[1]
Define the forward difference operator.
\end{exercise}

\begin{exercise}[1]
Determine the order of

\[
x_{n+3}
-
2x_{n+1}
+
x_n
=
0.
\]
\end{exercise}

\begin{exercise}[1]
Determine whether

\[
x_{n+1}
=
3x_n+2
\]

is linear or nonlinear.
\end{exercise}

\begin{exercise}[1]
Determine whether

\[
x_{n+1}=x_n^2+1
\]

is linear or nonlinear.
\end{exercise}

\begin{exercise}[1]
Define a fixed point of a recurrence.
\end{exercise}

\begin{exercise}[1]
State the stability criterion for a scalar fixed point.
\end{exercise}

\begin{exercise}[1]
Define a characteristic equation.
\end{exercise}

\begin{exercise}[1]
Define a generating function.
\end{exercise}

\begin{exercise}[1]
Define the one-sided \(z\)-transform.
\end{exercise}

\begin{exercise}[2]
Solve

\[
x_{n+1}=4x_n,
\qquad
x_0=3.
\]
\end{exercise}

\begin{exercise}[2]
Solve

\[
x_{n+1}=x_n+5,
\qquad
x_0=2.
\]
\end{exercise}

\begin{exercise}[2]
Find the equilibrium of

\[
x_{n+1}
=
0.8x_n+10.
\]
\end{exercise}

\begin{exercise}[2]
Determine whether the previous equilibrium is stable.
\end{exercise}

\begin{exercise}[2]
Solve

\[
x_{n+1}
=
\frac13x_n+2,
\qquad
x_0=0.
\]
\end{exercise}

\begin{exercise}[2]
Find the characteristic equation for

\[
x_{n+2}
-
7x_{n+1}
+
12x_n
=
0.
\]
\end{exercise}

\begin{exercise}[2]
Solve the previous recurrence.
\end{exercise}

\begin{exercise}[2]
Solve

\[
x_{n+2}
-
6x_{n+1}
+
9x_n
=
0.
\]
\end{exercise}

\begin{exercise}[2]
Find the fixed points of

\[
x_{n+1}
=
2x_n(1-x_n).
\]
\end{exercise}

\begin{exercise}[2]
Determine the stability of the fixed points in the previous problem.
\end{exercise}

\begin{exercise}[3]
Derive

\[
x_n
=
a^nx_0
+
b
\frac{
1-a^n
}{
1-a
}
\]

for

\[
x_{n+1}=ax_n+b.
\]
\end{exercise}

\begin{exercise}[3]
Solve

\[
x_{n+2}
-
5x_{n+1}
+
6x_n
=
0,
\]

with

\[
x_0=2,
\qquad
x_1=7.
\]
\end{exercise}

\begin{exercise}[3]
Solve

\[
x_{n+2}
+
4x_n
=
0.
\]
\end{exercise}

\begin{exercise}[3]
Solve

\[
x_{n+2}
-
4x_{n+1}
+
4x_n
=
0,
\]

with specified initial data of your choice.
\end{exercise}

\begin{exercise}[3]
Solve

\[
x_{n+1}
-
3x_n
=
2.
\]
\end{exercise}

\begin{exercise}[3]
Solve

\[
x_{n+1}
-
2x_n
=
3\cdot2^n
\]

using a resonant particular solution.
\end{exercise}

\begin{exercise}[3]
Derive Binet's formula for Fibonacci numbers.
\end{exercise}

\begin{exercise}[3]
Derive the Fibonacci generating function

\[
G(z)
=
\frac{z}{1-z-z^2}.
\]
\end{exercise}

\begin{exercise}[3]
Derive the one-sided \(z\)-transform shift formula

\[
\mathcal{Z}\{x_{n+1}\}
=
z[X(z)-x_0].
\]
\end{exercise}

\begin{exercise}[3]
Solve a first-order recurrence using the \(z\)-transform.
\end{exercise}

\begin{exercise}[3]
Solve

\[
\mathbf{x}_{n+1}
=
\begin{pmatrix}
2&0\\
0&1/2
\end{pmatrix}
\mathbf{x}_n.
\]
\end{exercise}

\begin{exercise}[3]
Classify the stability of the origin in the previous system.
\end{exercise}

\begin{exercise}[3]
Derive the discrete variation-of-constants formula.
\end{exercise}

\begin{exercise}[3]
Find the equilibrium of

\[
\mathbf{x}_{n+1}
=
A\mathbf{x}_n+\mathbf{b}.
\]
\end{exercise}

\begin{exercise}[3]
Derive the future-value formula for regular periodic deposits.
\end{exercise}

\begin{exercise}[3]
Formulate a monthly loan-balance recurrence.
\end{exercise}

\begin{exercise}[3]
Derive the fixed points of the classical logistic map.
\end{exercise}

\begin{exercise}[3]
Show that the nonzero logistic-map fixed point is stable for

\[
1<R<3.
\]
\end{exercise}

\begin{exercise}[3]
Explain the meaning of a period-two orbit.
\end{exercise}

\begin{exercise}[3]
Derive the stability multiplier for a period-\(k\) orbit.
\end{exercise}

\begin{exercise}[3]
Derive the Euler difference equation for

\[
y'=-ky.
\]
\end{exercise}

\begin{exercise}[3]
Show that Euler's method for

\[
y'=-ky
\]

is stable only when

\[
0<hk<2.
\]
\end{exercise}

\begin{exercise}[3]
Derive the centered finite-difference approximation for \(f''(x)\).
\end{exercise}

\begin{exercise}[4]
Study asymptotic theory of linear recurrences with several dominant
roots.
\end{exercise}

\begin{exercise}[4]
Study recurrences using companion matrices.
\end{exercise}

\begin{exercise}[4]
Derive Jordan-form solutions of matrix difference equations.
\end{exercise}

\begin{exercise}[4]
Study Perron--Frobenius theory for nonnegative recurrence matrices.
\end{exercise}

\begin{exercise}[4]
Analyze Leslie matrix population models.
\end{exercise}

\begin{exercise}[4]
Study stationary distributions of finite Markov chains.
\end{exercise}

\begin{exercise}[4]
Investigate convergence rates using subdominant eigenvalues.
\end{exercise}

\begin{exercise}[4]
Develop discrete Green's functions for linear recurrences.
\end{exercise}

\begin{exercise}[4]
Study discrete convolution and transfer functions.
\end{exercise}

\begin{exercise}[4]
Investigate regions of convergence for bilateral \(z\)-transforms.
\end{exercise}

\begin{exercise}[4]
Study pole-zero analysis of discrete linear systems.
\end{exercise}

\begin{exercise}[4]
Analyze period-doubling bifurcations in nonlinear maps.
\end{exercise}

\begin{exercise}[4]
Construct a logistic-map bifurcation diagram.
\end{exercise}

\begin{exercise}[4]
Compute Lyapunov exponents numerically for discrete maps.
\end{exercise}

\begin{exercise}[4]
Study the relation between Euler discretization and dynamical-system
stability.
\end{exercise}

\begin{exercise}[4]
Analyze higher-order finite-difference approximations.
\end{exercise}

\begin{exercise}[4]
Study consistency, stability, and convergence of difference schemes.
\end{exercise}

%%%%%%%%%%%%%%%%%%%%%%%%%%%%%%%%%%%%%%%%%%%%%%%%%%%%%%%%%%%%
\subsection{Conceptual Summary}
\label{subsec:difference-equation-summary}
%%%%%%%%%%%%%%%%%%%%%%%%%%%%%%%%%%%%%%%%%%%%%%%%%%%%%%%%%%%%

A difference equation governs a discrete sequence:

\[
\boxed{
x_{n+1}
=
f(x_n).
}
\]

The forward difference is

\[
\boxed{
\Delta x_n
=
x_{n+1}-x_n.
}
\]

A first-order affine recurrence

\[
\boxed{
x_{n+1}=ax_n+b
}
\]

has equilibrium

\[
\boxed{
x^*
=
\frac{
b
}{
1-a
},
\qquad
a\neq1,
}
\]

and solution

\[
\boxed{
x_n
=
x^*
+
a^n(x_0-x^*).
}
\]

A constant-coefficient homogeneous recurrence is solved using the trial

\[
\boxed{
x_n=r^n.
}
\]

This produces a characteristic polynomial.

For a second-order equation with distinct roots,

\[
\boxed{
x_n
=
C_1r_1^n
+
C_2r_2^n.
}
\]

Repeated roots require factors of \(n\).

Complex roots produce discrete oscillation.

A scalar fixed point satisfies

\[
\boxed{
f(x^*)=x^*,
}
\]

and is locally attracting when

\[
\boxed{
|f'(x^*)|<1.
}
\]

A linear vector recurrence is

\[
\boxed{
\mathbf{x}_{n+1}
=
A\mathbf{x}_n,
}
\]

with solution

\[
\boxed{
\mathbf{x}_n
=
A^n\mathbf{x}_0.
}
\]

Its origin is asymptotically stable when

\[
\boxed{
\rho(A)<1.
}
\]

Difference equations therefore combine

\[
\boxed{
\text{sequences}
+
\text{linear algebra}
+
\text{discrete dynamics}
+
\text{transforms}
+
\text{numerical methods}.
}
\]

%%%%%%%%%%%%%%%%%%%%%%%%%%%%%%%%%%%%%%%%%%%%%%%%%%%%%%%%%%%%
\subsection{Section Connections}
\label{subsec:difference-equation-connections}
%%%%%%%%%%%%%%%%%%%%%%%%%%%%%%%%%%%%%%%%%%%%%%%%%%%%%%%%%%%%

\chapterconnection{
Difference Equations provide the discrete-time counterpart of ordinary
differential equations. Characteristic roots parallel characteristic
exponents, matrix powers parallel matrix exponentials, and the unit
circle replaces the left half-plane as the natural stability boundary.
Generating functions and the \(z\)-transform provide transform methods
for recurrences, while nonlinear maps connect directly with the
dynamical-systems ideas of Section~\ref{sec:dynamical-systems}. Difference
equations also arise naturally when differential equations are
discretized numerically. The next section, Integral Equations, studies
equations in which the unknown function appears under an integral and
develops connections with differential equations, Green's functions,
operator theory, and transform methods.
}

%%%%%%%%%%%%%%%%%%%%%%%%%%%%%%%%%%%%%%%%%%%%%%%%%%%%%%%%%%%%
% SECTION SOURCE TRACKING
%%%%%%%%%%%%%%%%%%%%%%%%%%%%%%%%%%%%%%%%%%%%%%%%%%%%%%%%%%%%

%------------------------------------------------------------
% 9.5 Difference Equations
%------------------------------------------------------------
%
% Source basis:
%   - Standard recurrence and difference-equation theory.
%   - Standard discrete dynamical-systems theory.
%   - Standard generating-function and z-transform methods.
%
% Used for:
%   - difference-equation definitions;
%   - forward and higher differences;
%   - order and initial conditions;
%   - linear and nonlinear recurrences;
%   - autonomous and nonautonomous equations;
%   - first-order linear recurrences;
%   - geometric sequences;
%   - affine recurrences;
%   - equilibria and fixed points;
%   - fixed-point stability;
%   - constant-coefficient recurrences;
%   - characteristic equations;
%   - distinct, repeated, and complex roots;
%   - superposition;
%   - nonhomogeneous recurrences;
%   - undetermined coefficients;
%   - discrete resonance;
%   - variable-coefficient first-order recurrences;
%   - telescoping sums;
%   - Fibonacci recurrence and Binet formula;
%   - generating functions;
%   - z-transform methods;
%   - discrete transfer functions;
%   - matrix difference equations;
%   - spectral-radius stability;
%   - discrete variation of constants;
%   - population and financial models;
%   - logistic maps;
%   - periodic orbits;
%   - period doubling;
%   - Euler discretization;
%   - finite differences;
%   - Markov-type recurrences;
%   - stationary distributions;
%   - Leslie matrices;
%   - discrete convolution;
%   - applications and exercises.
%
% Notes:
%   - Integral Equations is developed next in Section 9.6.
%   - Delay Differential Equations and Differential-Algebraic Equations
%     follow later in Chapter 9.
%   - Numerical difference schemes are developed more deeply in
%     Chapter 11.
%
%%%%%%%%%%%%%%%%%%%%%%%%%%%%%%%%%%%%%%%%%%%%%%%%%%%%%%%%%%%%